\section{Conventions and analytic inputs}\label{sec:inputs}

Throughout the proof $x$ is a real parameter tending to infinity, and
\[
 L=\log x,\qquad W=\exp(L^{0.24}),\qquad
 V(W)=\prod_{p\le W}\left(1-\frac1p\right),\qquad
 \e(t)=\exp(2\pi i t).
\]
A positive integer is \emph{rough} if none of its prime factors is at
most $W$; we set $P^-(1)=\infty$.  Prime variables in explicitly
indicated prime sums range only over primes.  The functions $\tau(n)$
and $\tau_j(n)$ count divisors and ordered factorizations into $j$
positive integers, respectively.  Dirichlet characters are zero on
nonunits.  We use $n\asymp H$ for a fixed bounded enlargement of
$[H,2H]$.

Every unspecified exponent of $L$ is fixed before $x$ tends to
infinity.  Constants may depend on previously fixed data.  Whenever
an exponent must be independent of a marking parameter, this is
stated explicitly.  The number of determinant pads in
Sections~\ref{sec:determinant}--\ref{sec:major-arcs} grows slowly
with $L$; the padding parameter in Section~\ref{sec:correlations}
is instead a sufficiently large fixed integer.  These are separate
constructions.

We use the dilation-graph and ideal-kernel theorems of the companion
\emph{Weighted dilation graphs, smooth shifted primes and totient
fibers}~\cite{SmoothShifted}, denoted S. Its complete text accompanies
this article. Section~\ref{sec:correlations} states the exact operator
contracts and verifies their endpoint hypotheses. The long-prime and
logarithmic-phase estimates stated below are proved in
Appendix~\ref{lp:section}. Our presieving cutoff is always the $W$
defined above; it is not S's cutoff $\exp(\sqrt{\log x})$.

\subsection{Prime estimates and mean squares}

\begin{lemma}[Prime estimates]\label{lem:prime-estimates}
The prime number theorem holds with an error smaller than every
fixed negative power of the logarithm.  In particular, as $y\to\infty$,
\[
 \sum_{p\le y}\frac1p=\log\log y+c_M+o(1),
 \qquad
 \prod_{p\le y}\left(1-\frac1p\right)
       \sim\frac{e^{-\gamma_E}}{\log y},
\]
where $c_M$ and $\gamma_E$ are constants.  For fixed $A,C>0$,
uniformly for $r\le(\log y)^C$ and $(a,r)=1$,
\[
 \#\{p\le y:p\equiv a\pmod r\}
 =\frac{\Li(y)}{\varphi(r)}
     +O_{A,C}\!\left(\frac{y}{(\log y)^A}\right).
\]
The constants need not be effective.
\end{lemma}

These classical estimates follow from the quantitative prime number
theorem, Siegel--Walfisz, and Mertens' theorems; see
\cite[Corollary~39 and Exercises~40, 64]{TaoSW} and
\cite[Theorems~15, 26]{TaoMertens}. The harmonic asymptotic also follows
by partial summation of the prime number theorem. Thus, for every fixed $a>0$,
\begin{equation}\label{eq:band-mass}
 \sum_{\exp(L^a)\le p\le\exp(2L^a)}\frac1p
       =\log 2+o(1).
\end{equation}
We use prime asymptotics only on fixed positive-power ranges or on
intervals for which taking differences of the stated estimates gives
the required absolute error.

For the classical, stronger coefficient-independent mean-value
theorem, see \cite[Theorem~2 and Corollary~3]{MontgomeryVaughan1974}.
We prove the weaker bounds needed here directly.

\begin{lemma}[Divisor moments and Dirichlet mean squares]
\label{lem:moment-bounds}
For every fixed nonnegative integer $k$ there is $C_k$ such that
\[
 \sum_{n\le Y}\tau(n)^k\ll_k Y(\log(2Y))^{C_k},\qquad
 \sum_{n\le Y}\frac{\tau(n)^k}{n}
       \ll_k(\log(2Y))^{C_k}.
\]
Let $P(t)=\sum_{n\le Y}c_n n^{it}$.  On every real interval $I$
of length $T\ge0$,
\begin{align}
 \int_I|P(t)|^2\dd t
 &\ll (T+Y\log(2Y))\sum_n|c_n|^2,\label{eq:dirichlet-mean}\\
 \sum_{t\in\mathcal T}|P(t)|^2
 &\ll (T+1+Y\log(2Y))(\log(2Y))^2
                  \sum_n|c_n|^2,\label{eq:dirichlet-spaced}
\end{align}
whenever $\mathcal T\subset I$ has mutual spacing at least one.
The constants in these two inequalities are absolute and do not
depend on how the coefficients were formed.
\end{lemma}

\begin{proof}
Unique factorization and positivity bound the harmonic divisor sum by
\[
 \prod_{p\le Y}\sum_{a\ge0}\frac{(a+1)^k}{p^a}
 =\prod_{p\le Y}\left(1+\frac{2^k}{p}+O_k(p^{-2})\right)
 \ll_k (\log(2Y))^{C_k}.
\]
Multiplication by $Y$ gives the counting bound. For the continuous
mean square, expand the square and integrate. The diagonal is
$T\sum_n|c_n|^2$. For $m\ne n$, the integral has absolute value at
most $2/|\log(n/m)|\ll Y/|n-m|$. The inequality
$2|c_nc_m|\le |c_n|^2+|c_m|^2$ and the harmonic sum over $|n-m|$
give \eqref{eq:dirichlet-mean}. On the unit interval centered at each
point of $\mathcal T$, the one-dimensional Sobolev inequality bounds
$|P(t)|^2$ by a constant times the integral of $|P|^2+|P'|^2$.
These intervals have bounded overlap and lie in an interval of length
$T+1$. Apply the continuous estimate to $P$ and $P'$, whose
coefficients are $i(\log n)c_n$, to obtain
\eqref{eq:dirichlet-spaced}.
\end{proof}

The coefficient-independent formulation in
\eqref{eq:dirichlet-mean}--\eqref{eq:dirichlet-spaced} is important:
later a small-prime polynomial is raised to a growing power, whose
coefficient norm is bounded separately by a factorial estimate.
Fixed divisor moments alone are not applied at that growing order.

\subsection{Prime polynomials and logarithmic phases}

\begin{lemma}[Long prime polynomial]\label{lem:long-prime}
Fix $0<\tau<\eta<1$ and $C,A>0$.  There is
$B_0=B_0(\tau,\eta,C,A)$ such that, uniformly for
\[
 \frac{x^\tau}{2}\le N\le x^\eta,\qquad
 q\le L^C,\qquad L^{B_0}\le |t|\le x^2,
\]
every Dirichlet character $\chi$ modulo $q$ and every interval
$I\subset[N,2N]$ satisfy
\begin{equation}\label{eq:long-prime}
 \left|\sum_{p\in I}\chi(p)p^{-1+it}\right|
       \ll_{\tau,\eta,C,A}L^{-A}.
\end{equation}
\end{lemma}

The complete proof is Lemma~\ref{lp:prime-polynomial} in
Appendix~\ref{lp:section}, with exactly the displayed ranges. Partial
summation converts the reciprocal weight in \eqref{eq:long-prime}
to normalization by $N$ on a dyadic interval.

\begin{lemma}[Logarithmic phases on progressions]\label{lem:log-phase}
Fix $C>0$.  There is an absolute constant $C_3>0$ such that, for
sufficiently large $x$ in terms of $C$, the following holds uniformly:
\[
 \exp\!\left(\frac{L}{(\log L)^2}\right)\le N\le2x^5,
 \quad q\le L^C,
 \quad
 \exp\!\left(\frac{L}{2(\log L)^2}\right)\le |t|\le4x^3.
\]
For every residue $a\pmod q$ and every interval $I\subset[N,2N]$,
\begin{equation}\label{eq:log-phase}
 \left|\sum_{\substack{n\in I\\n\equiv a\pmod q}}n^{it}\right|
 \ll \frac Nq\exp\!\left(-\frac{L}{(\log L)^{C_3}}\right).
\end{equation}
\end{lemma}

The proof is Lemma~\ref{lp:log-phase} in Appendix~\ref{lp:section}.
The estimate applies to arbitrary residue
classes, not only units.  At lower frequencies we use the elementary
sum--integral comparison on a progression.  For $N\ge x^\delta$,
$q\le L^C$ and $1\le|t|\le\exp(L/(\log L)^2)$, it gives
\begin{equation}\label{eq:moderate-phase}
 \frac qN\left|\sum_{\substack{n\in I\\n\equiv a\pmod q}}n^{it}\right|
 \ll \frac1{|t|}+x^{-c_\delta},
\end{equation}
after decreasing $c_\delta>0$ if necessary.  Indeed the integral is
$O(N/(q|t|))$, and the endpoint/total-variation error is
$O(1+|t|)$.  The same estimate holds after partial summation against
a smooth factor with fixed logarithmic derivative costs, with those
costs included in the implied logarithmic power.

\subsection{Elementary rough counts}

\begin{lemma}[Rough integers in long intervals]\label{lem:rough-counts}
Fix $\delta>0$, $C>0$, $0<a<1$ and $A>0$.  Suppose
$x^\delta\le H\le x^C$ and $I\subset[H,2H]$ is an interval.
Then
\begin{equation}\label{eq:rough-count}
 \#\{n\in I:P^-(n)>W\}=V(W)|I|+O(HL^{-A}).
\end{equation}
If $d\le\exp(L^a)$ and every prime divisor of $d$ exceeds $W$,
then for every residue $b\pmod d$,
\begin{equation}\label{eq:rough-off-sieve-progression}
 \#\{n\in I:n\equiv b\pmod d,\ P^-(n)>W\}
 =\frac{V(W)|I|}{d}+O\!\left(\frac Hd L^{-A}\right).
\end{equation}
For every character $\chi$ modulo $q\le L^C$,
\begin{equation}\label{eq:rough-character}
 \sum_{\substack{n\in I\\P^-(n)>W}}\chi(n)
 =\1_{\chi\text{ principal}}V(W)|I|+O(HL^{-A}).
\end{equation}
All statements are uniform in the indicated intervals and residues.
\end{lemma}

\begin{proof}
Use consecutive even and odd Bonferroni truncations for the events
$p\mid n$, $p\le W$, at depths $h$ and $h+1$, with
$h\asymp C_A\log L$.  Their divisors are at most
\[
 W^{h+1}=\exp\bigl(O_A(L^{0.24}\log L)\bigr)=x^{o(1)}.
\]
The number of divisor terms and their total absolute integer
coefficient mass are also $x^{o(1)}$.  Counting each progression
by its length divided by the modulus costs at most one per term.
Writing $S_W=\sum_{p\le W}1/p=O(\log L)$, the difference of the
two model bounds is at most
\[
 \frac{S_W^{h+1}}{(h+1)!},
\]
which is smaller than any prescribed power of $L^{-1}$ on choosing
the fixed constant $C_A$ sufficiently large.  The full model product
is $V(W)$.

For \eqref{eq:rough-off-sieve-progression}, each sieve divisor is
coprime to $d$, so the Chinese remainder theorem gives main term
$|I|/(de)$ for the divisor $e$.  The same omitted-degree bound is
multiplied by $H/d$.  The total endpoint error remains $x^{o(1)}$,
which is smaller than $(H/d)L^{-A}$, since $d=x^{o(1)}$ and
$H\ge x^\delta$.  This proves both counting formulas.

For \eqref{eq:rough-character}, first count a fixed unit residue
class modulo $q$.  Primes dividing $q$ are automatically absent;
the model density in that progression is
\[
 \frac1q\prod_{\substack{p\le W\\p\nmid q}}
        \left(1-\frac1p\right)=\frac{V(W)}{\varphi(q)}.
\]
The endpoint and Bonferroni errors have arbitrary logarithmic
precision, uniformly for $q\le L^C$.  Sum with the character over
the unit residue classes, choosing the precision first to absorb
their number.  Character orthogonality gives the displayed answer.
\end{proof}

\subsection{A sieve for nonnegative weighted objects}

\begin{lemma}[Block sieve]\label{lem:block-sieve}
Consider finitely many objects with nonnegative weights and a bad
condition at each of some designated primes $p\le z$.  Suppose
that the total weight on which all conditions indexed by the primes
dividing a squarefree $d$ hold is
\[
 Xg(d)+r(d),
\]
including $d=1$, where $X\ge0$, $g$ is multiplicative, and fixed
$\eta_0>0,C_0$ satisfy
\[
 0\le g(p)\le1-\eta_0,
 \qquad \sum_{w<p\le w^2}g(p)\le C_0\quad(w>1).
\]
For every sufficiently large even integer $h$, the weight avoiding
all designated bad conditions is
\begin{equation}\label{eq:block-sieve}
 X\prod_{p\le z}(1-g(p))\bigl(1+O(e^{-h})\bigr)
 +O\!\left(\sum_{\substack{d\le z^{4h+2}\\d\text{ squarefree}}}
          |r(d)|\right).
\end{equation}
Products and sums use only designated primes.  Constants depend
only on the density bounds and are uniform in $h$.  For $h=2$,
an upper bound holds with a fixed constant times the main term
and the same remainder sum.
\end{lemma}

This is the block sieve of S~\cite[Lemma~2.9]{SmoothShifted}.
Its upper and lower polynomials have coefficients of absolute value
at most one, supported on squarefree $d\le z^{4h+2}$; the upper
polynomial is nonnegative.  We will supply all needed remainder
estimates for our particular weights.  No assertion detecting
primes from congruence information alone is included in this input.
